\chapter{Solutions: Midterm Exam, April 30, 2015}
\label{s:physics-2015-04-midterm}
These solutions are derived from the stated assumptions. The density of states $D(E)$ counts states per unit energy in the entire volume $V$. The factor $g_s$ denotes any spin degeneracy included separately; the free-electron gas uses $g_s=2$ without double counting.

\section*{1. Lattice translations in BaTiO$_3$}
For the cubic model specified in the question, the Bravais translations are
$\vct R=a(n_x,n_y,n_z)$ with $n_i\in\mathbb Z$. A five-atom basis is
\[
\begin{array}{c|c}
\text{Atom}&\vct\tau/a\\ \hline
\mathrm{Ba}&(0,0,0)\\
\mathrm{Ti}&(1/2,1/2,1/2)\\
\mathrm{O}_1&(1/2,1/2,0)\\
\mathrm{O}_2&(1/2,0,1/2)\\
\mathrm{O}_3&(0,1/2,1/2)
\end{array}
\]
Actual atomic positions are $\vct R+\vct\tau_\nu$. A displacement is a Bravais translation precisely when it belongs to $a\mathbb Z^3$.
\begin{enumerate}[label=(\alph*)]
\item \textbf{True.} There is one Ti basis position in each primitive cell, so subtraction cancels the basis position:
\[
(\vct R_m+\vct\tau_{\rm Ti})-(\vct R_n+\vct\tau_{\rm Ti})=\vct R_m-\vct R_n.
\]
This is always a lattice translation.
\item \textbf{True.} A Ba-to-Ti displacement has the form
$a(n_x+1/2,n_y+1/2,n_z+1/2)$. Its half-integer components exclude it from $a\mathbb Z^3$.
\item \textbf{False.} Two O atoms with the same basis index are related by a translation, but the assertion concerns any two O atoms. Within one cell,
\[
\vct\tau_{\mathrm O_2}-\vct\tau_{\mathrm O_1}=\frac a2(0,-1,1),
\]
which is not a simple-cubic lattice translation. Having the same chemical species does not imply having the same basis index.
\end{enumerate}

\section*{2. The fcc lattice, its reciprocal lattice, and lattice planes}
\subsection*{(a) Primitive-cell volumes}
The given vectors are primitive vectors of a face-centered cubic (fcc) lattice. Their scalar triple product is
\[
\Omega=\abs{\vct t_1\cdot(\vct t_2\times\vct t_3)}
=\left(\frac a2\right)^3
\det\begin{pmatrix}0&1&1\\1&0&1\\1&1&0\end{pmatrix}
=\frac{a^3}{4}.
\]
A Wigner--Seitz cell is also a primitive cell. The reciprocal primitive-cell volume, equal to the first Brillouin-zone volume, is $(2\pi)^3/\Omega$. Thus
\[
\boxed{V_{\rm WS}=\frac{a^3}{4}},\qquad
\boxed{V_{\rm BZ}=\frac{32\pi^3}{a^3}}.
\]

\subsection*{(b) Reciprocal primitive vectors}
Using $\vct g_1=2\pi(\vct t_2\times\vct t_3)/\Omega$ and cyclic permutations gives
\[
\boxed{\vct g_1=\frac{2\pi}{a}(-1,1,1),\quad
\vct g_2=\frac{2\pi}{a}(1,-1,1),\quad
\vct g_3=\frac{2\pi}{a}(1,1,-1)}.
\]
For example, $\vct t_1\cdot\vct g_1=2\pi$ and
$\vct t_1\cdot\vct g_2=\vct t_1\cdot\vct g_3=0$. The other rows give the same reciprocal orthogonality relation.

\subsection*{(c) Spacing of successive occupied lattice planes}
The unit normal is $\hat{\vct n}=(0,1,1)/\sqrt2$. Project an arbitrary lattice point
$\vct R=n_1\vct t_1+n_2\vct t_2+n_3\vct t_3$ onto this direction:
\[
\hat{\vct n}\cdot\vct R
=\frac{a}{2\sqrt2}(2n_1+n_2+n_3).
\]
The integer in parentheses can take every integer value. Successive planes containing fcc lattice points therefore lie at
$\hat{\vct n}\cdot\vct R=j a/(2\sqrt2)$, giving
\[
\boxed{d=\frac{a}{2\sqrt2}}.
\]
Equivalently, the shortest nonzero reciprocal vector parallel to the normal is
$\vct G=2\vct g_1+\vct g_2+\vct g_3=(4\pi/a)(0,1,1)$, and $d=2\pi/|\vct G|$.
\begin{sourceNote}
Here ``successive lattice planes'' means successive occupied fcc layers. The formal conventional-cubic $(011)$ spacing is $a/\sqrt2$, but fcc centering supplies an additional occupied layer halfway between those planes. The occupied-layer spacing requested under this interpretation is consequently half as large.
\end{sourceNote}

\subsection*{(d) Areal density}
Two perpendicular in-plane translations are
\[
\vct u=a(1,0,0)=-\vct t_1+\vct t_2+\vct t_3,\qquad
\vct v=\frac a2(0,1,-1)=\vct t_3-\vct t_2.
\]
They generate the two-dimensional lattice on this plane. Its primitive area is
$A_{\parallel}=|\vct u\times\vct v|=a^2/\sqrt2$, with one lattice point per primitive parallelogram. Hence
\[
\boxed{\sigma=\frac1{A_{\parallel}}=\frac{\sqrt2}{a^2}}.
\]
The volume check is $A_{\parallel}d=\Omega$, equivalently $\sigma=d/\Omega$.
\begin{center}
\begin{tikzpicture}[>=Latex]
\draw[thick] (0,0) rectangle (3,2.12);
\foreach \x in {0,3}\foreach \y in {0,2.12}\fill (\x,\y) circle (2pt);
\draw[->,blue!60!black,thick] (0,0)--(3,0) node[midway,below] {$\vct u,\ |\vct u|=a$};
\draw[->,blue!60!black,thick] (0,0)--(0,2.12) node[midway,left] {$\vct v$};
\node[right] at (3,1.06) {$|\vct v|=a/\sqrt2$};
\node at (1.5,1.06) {$A_{\parallel}=a^2/\sqrt2$};
\node[below] at (1.5,-.5) {Four corners contribute $1/4$ point each.};
\end{tikzpicture}
\end{center}

\subsection*{(e) Testing the proposed vectors}
Every reciprocal vector must satisfy
$\vct t_i\cdot\vct G\in2\pi\mathbb Z$ for all primitive translations. However,
\[
\vct t_1\cdot\vct G_1
=\frac a2(0,1,1)\cdot\frac\pi a(0,1,1)=\pi.
\]
Thus $\vct G_1$ is not even a reciprocal-lattice vector. The proposed vectors \textbf{do not} generate the required reciprocal lattice by integer linear combinations. Linear independence over the real numbers only guarantees that three vectors span three-dimensional vector space; it does not establish equality of discrete lattices.

\section*{3. Atomic density of diamond}
Diamond consists of an fcc Bravais lattice with a two-atom basis. A conventional cubic cell contains
$8(1/8)+6(1/2)=4$ fcc lattice points. With two atoms per point,
\[
\boxed{n_{\rm atom}=\frac{4\times2}{a^3}=\frac8{a^3}}.
\]
The equivalent primitive-cell count is two atoms divided by $a^3/4$.

\section*{4. Definitions}
\begin{enumerate}[label=(\alph*)]
\item \textbf{Three-dimensional hexagonal lattice.} One choice of primitive vectors is
\[
\vct a_1=a(1,0,0),\quad
\vct a_2=a(-1/2,\sqrt3/2,0),\quad
\vct a_3=c(0,0,1).
\]
The first two vectors have equal lengths and enclose $120^\circ$; the third is perpendicular to their plane. Their integer combinations define the hexagonal Bravais lattice, of primitive volume $\sqrt3a^2c/2$. Hexagonal close packing (hcp) adds a two-atom basis to this lattice; the lattice and its basis are separate concepts.
\item \textbf{Bloch theorem.} If a single-electron Hamiltonian is invariant under every lattice translation, its energy eigenstates can be chosen as
\[
\psi_{n\vct k}(\vct r)=\ee^{\ii\vct k\cdot\vct r}u_{n\vct k}(\vct r),
\qquad u_{n\vct k}(\vct r+\vct R)=u_{n\vct k}(\vct r).
\]
Equivalently,
$\psi_{n\vct k}(\vct r+\vct R)=\ee^{\ii\vct k\cdot\vct R}\psi_{n\vct k}(\vct r)$.
The wave vector labels crystal momentum and can be restricted to the first Brillouin zone modulo a reciprocal vector.
\item \textbf{Critical point.} A point where $\nabla_{\vct k}\varepsilon_n=0$. For a nondegenerate Hessian in three dimensions, the type $M_j$ counts its negative eigenvalues: $M_0$ is a minimum, $M_1$ and $M_2$ are saddles with one and two negative directions, and $M_3$ is a maximum.
\item \textbf{Fermi energy.} The highest occupied single-particle energy after filling the zero-temperature ground state at fixed particle number. For the free-electron gas here it equals $\mu(0)=\varepsilon_F$. Its energy reference must be specified; it is distinct from the work function required to remove an electron to vacuum.
\end{enumerate}

\section*{5. Density of states near a three-dimensional $M_1$ saddle}
\subsection*{(a) Local dispersion}
Set $\vct q=\vct k-\vct k_1$ and diagonalize the Hessian. The gradient vanishes at the critical point, so no linear term remains. With one negative-curvature direction,
\[
\boxed{\varepsilon(\vct k)
=\varepsilon_1+Aq_x^2+Bq_y^2-Cq_z^2+O(q^3)},
\qquad A,B,C>0.
\]
The coefficients $A,B,C$ already incorporate the factor $1/2$ in the Taylor expansion.

\subsection*{(b) Deriving the nonanalytic contribution with a local cutoff}
Define $X=\sqrt A q_x$, $Y=\sqrt B q_y$, $Z=\sqrt C q_z$, and $\Delta=E-\varepsilon_1$. Then
\[
D_{\rm loc}(E)=\frac{g_sV}{(2\pi)^3\sqrt{ABC}}
\int\dd X\dd Y\dd Z\,
\delta(\Delta-X^2-Y^2+Z^2).
\]
The quadratic saddle cannot be extended to infinity: its constant-energy surfaces are unbounded. Restrict it to a neighborhood, for example $|Z|<\Lambda$, with a transverse cutoff large enough to contain the relevant energy shells in the range considered. Polar integration over $X,Y$ gives
\[
\int_0^{\infty}2\pi\rho\dd\rho\,
\delta(\Delta+Z^2-\rho^2)=\pi\Theta(\Delta+Z^2).
\]
For $\Delta>0$, every $Z$ in $(-\Lambda,\Lambda)$ contributes. For
$-\Lambda^2<\Delta<0$, the interval $|Z|<\sqrt{-\Delta}$ has no real transverse radius and must be removed. Therefore
\[
D_{\rm loc}(E)=\frac{g_sV}{4\pi^2\sqrt{ABC}}
\left[\Lambda-\sqrt{\varepsilon_1-E}\,
\Theta(\varepsilon_1-E)\right].
\]
Other cutoff choices and other parts of the band change the regular background, while the leading local nonanalytic term is
\[
\boxed{D(E)=D_{\rm reg}(E)
-\frac{g_sV}{4\pi^2\sqrt{ABC}}
\sqrt{\varepsilon_1-E}\,\Theta(\varepsilon_1-E)}.
\]
This expression gives the leading behavior near the critical energy; higher-order terms in the dispersion can add subleading corrections. The density of states remains finite at $E=\varepsilon_1$, but its slope diverges on the lower-energy side. Local curvature alone cannot fix the absolute background $D_{\rm reg}$ of the whole band.
\begin{center}
\begin{tikzpicture}[>=Latex]
\draw[->] (0,0)--(5,0) node[right] {$E$};
\draw[->] (0,0)--(0,3) node[above] {$D_{\rm loc}$};
\draw[dashed,gray] (3,0)--(3,2.4);
\draw[thick,blue!65!black,domain=.35:3,samples=100] plot (\x,{2.4-sqrt(3-\x)});
\draw[thick,blue!65!black] (3,2.4)--(4.8,2.4);
\node[below] at (3,0) {$\varepsilon_1$};
\node[above] at (2.7,2.55) {Finite value; divergent slope from below};
\end{tikzpicture}
\end{center}

\section*{6. Free-electron ground-state quantities in terms of $r_s$}
Let $a_0=4\pi\varepsilon_0\hbar^2/(m_e e^2)$ be the Bohr radius. By definition, a sphere of radius $r_sa_0$ contains the mean volume per electron:
\[
\frac1n=\frac{4\pi}{3}(r_sa_0)^3
\quad\Longrightarrow\quad
\boxed{n=\frac{3}{4\pi r_s^3a_0^3}}.
\]
Counting both spins inside the Fermi sphere gives
\[
N=2\frac V{(2\pi)^3}\frac{4\pi}{3}k_F^3,
\qquad n=\frac{k_F^3}{3\pi^2}.
\]
Consequently,
\begin{align*}
\boxed{k_F=\frac{(9\pi/4)^{1/3}}{r_sa_0}},\qquad
\boxed{\varepsilon_F=\frac{\hbar^2}{2m_ea_0^2}
\left(\frac{9\pi}{4}\right)^{2/3}\frac1{r_s^2}},\\
\boxed{T_F=\frac{\varepsilon_F}{\kb}}.
\end{align*}
The mean energy averages all occupied states, rather than only states on the Fermi surface:
\[
\overline\varepsilon
=\frac{\int_0^{k_F}(\hbar^2k^2/2m_e)k^2\dd k}
{\int_0^{k_F}k^2\dd k}
=\boxed{\frac35\varepsilon_F}.
\]

\section*{7. Low-temperature heat capacity and the printed coefficient}
For $T\ll T_F$, states well below the Fermi energy are almost completely filled. Pauli exclusion prevents arbitrary changes of their occupations. Thermal excitations predominantly involve an energy window of width of order $\kb T$ near $\varepsilon_F$.
Since $D(\varepsilon_F)=3N/(2\varepsilon_F)$ in three dimensions, the number of thermally active states scales as
\[
N_{\rm active}\sim D(\varepsilon_F)\kb T
=\frac32N\frac{T}{T_F}.
\]
Assigning a heat-capacity scale $\kb$ to each active state gives the estimate printed in the question: a suppression by $T/T_F$ relative to the classical value.

The printed factor $3/2$ is not the exact low-temperature coefficient. The Sommerfeld expansion at fixed $N,V$ gives
\[
C_V=\frac{\pi^2}{3}\kb^2T D(\varepsilon_F)
=\boxed{\frac{\pi^2}{2}N\kb\frac{T}{T_F}}
+O\!\left(N\kb(T/T_F)^3\right).
\]
Thus the physical argument establishes linear temperature dependence and a restricted thermally active population. It must not be interpreted as an exact derivation of the numerical coefficient $3/2$.

\section*{8. The exponential tail above the vacuum level}
The Fermi--Dirac distribution is
\[
f(E)=\frac1{\ee^{(E-\mu)/(\kb T)}+1}.
\]
Define the work function $\Phi=E_{\rm vac}-\mu>0$. When $\Phi\gg\kb T$, as in the usual low-temperature metallic regime, all $E\ge E_{\rm vac}$ satisfy
$x=(E-\mu)/(\kb T)\ge\Phi/(\kb T)\gg1$. It follows that
\[
f(E)=\frac{\ee^{-x}}{1+\ee^{-x}}
\simeq\boxed{\exp[-(E-\mu)/(\kb T)]}.
\]
The relative error of $f_{\rm exp}=\ee^{-x}$ is exactly
\[
\frac{f_{\rm exp}-f}{f}=\ee^{-x}
\le\ee^{-\Phi/(\kb T)}.
\]
The approximation is controlled by the energy separation from the chemical potential compared with thermal energy, not by the word ``vacuum'' alone.

\section*{References and checks}
The counting arguments and integrals are derived above. Primary teaching references for comparison are
\href{https://www.damtp.cam.ac.uk/user/tong/statphys/statmechhtml/S3.html}{David Tong, Statistical Physics, Chapter 3}
(state counting and the Sommerfeld expansion), and
\href{https://www.mit.edu/course/6/6.732/www/6.732-pt2.pdf}{M. S. Dresselhaus, MIT Solid State Physics Part II, Chapter 4}
(critical-point classification and the three-dimensional saddle-point singularity).
